\chapter{Wilson's Disease}
\index{Wilson's Disease}
\label{sec:Wilson's Disease}

\section{Overview}
Wilson's disease is an autosomal recessive disorder of copper metabolism caused by mutations in the ATP7B gene, with a prevalence of approximately 1 in 30,000.
\section{Causes}
This condition results from disruption of normal wilson's disease function.

\begin{itemize}[noitemsep]
  \item This condition happens because of impaired biliary copper excretion leading to accumulation of copper in the liver, brain, cornea, and other organs. liver manifestations range from asymptomatic transaminase elevation to acute hepatitis, chronic hepatitis, or cirrhosis
  \item Neuropsychiatric features include tremor, slurred speech, dystonia, parkinsonism, personality changes, and depression. Kayser-Fleischer rings (copper deposition in Descemet's membrane of the cornea) are a hallmark finding.
\end{itemize}
\section{Risk Factors}


\begin{itemize}[noitemsep]
  \item Genetic predisposition and family history
  \item Advancing age
  \item Lifestyle factors (diet, exercise, smoking, alcohol)
  \item Environmental exposures
  \item Underlying medical conditions
  \item Medication use
\end{itemize}
\section{Signs and Symptoms}

\subsection{Common symptoms}
\begin{itemize}[noitemsep]
  \item \item Pain or discomfort in the affected area    \item Pain or discomfort in the affected area
  \end{itemize}

\subsection{Emergency warning signs}
\begin{itemize}[noitemsep]
  \item Severe or worsening symptoms of wilson's disease require immediate medical evaluation
\end{itemize}
\section{Progression and Stages}
The condition may progress through different stages of severity over time. Early stages often involve mild or subtle symptoms that may be overlooked. Moderate stages involve more noticeable symptoms affecting daily function. Advanced stages may involve significant impairment and complications.
\section{Complications}
\begin{itemize}[noitemsep]
  \item Disease progression leading to worsening symptoms
  \item Secondary infections or injuries
  \item Side effects of treatment
  \item Reduced quality of life
  \item Emotional and psychological impact
\end{itemize}
\section{Diagnosis}
Diagnosis typically involves low serum ceruloplasmin (less than 20 mg/dL), elevated 24-hour urinary copper excretion, and liver copper quantification on liver biopsy. Comprehensive medical history and physical examination Laboratory tests to assess organ function and rule out other conditions Imaging studies to visualize affected structures Specialized testing depending on the affected system Consultation with relevant specialists Monitoring of disease progression over time

\begin{itemize}[noitemsep]
  \item Comprehensive medical history and physical examination
  \item Laboratory tests to assess organ function
  \item Imaging studies to visualize affected structures
  \item Specialized testing depending on the affected system
  \item Consultation with relevant specialists
\end{itemize}
\section{Prevention}
\begin{itemize}[noitemsep]
  \item Maintain a healthy lifestyle with balanced diet and exercise
  \item Avoid known risk factors
  \item Attend regular medical check-ups for early detection
  \item Follow recommended screening guidelines
\end{itemize}
\section{Treatment Options}
\begin{itemize}[noitemsep]
  \item Treatment options include copper chelation (penicillamine, trientine), zinc therapy to block copper absorption, and liver transplantation in fulminant or decompensated disease
  \item Screening of first-degree relatives is essential.
  \item Medications to manage symptoms and address underlying mechanisms
  \item Lifestyle modifications and self-care measures
  \item Physical therapy and rehabilitation
  \item Surgical intervention when indicated
  \item Regular monitoring and follow-up care
\end{itemize}
\section{Living with It}
\begin{itemize}[noitemsep]
  \item Follow your treatment plan as prescribed
  \item Maintain regular follow-up with your healthcare provider
  \item Make necessary lifestyle adjustments
  \item Seek support from family, friends, or support groups
  \item Stay informed about your condition
\end{itemize}
\section{Outlook}
Most people recover well with early diagnosis and lifelong treatment. The outlook varies depending on the specific condition, severity at diagnosis, and response to treatment. Early intervention generally leads to better outcomes. Many conditions can be effectively managed with appropriate treatment. Regular follow-up care is important for monitoring and treatment adjustment.
